% fig_tstar_sld.tex -- MKP, SLD: subgradient iterations t* needed to beat the GNN root bound, and time of these t* iterations
% vs. one GNN bound call. Medians over 10 instances per (n, N); generated from latex_data/time_vs_tstar_sld.csv
% (extract_time_vs_tstar.py). Preamble: \usepackage{pgfplots} \usepgfplotslibrary{groupplots} \pgfplotsset{compat=1.16}
% Usage in the paper: % fig_tstar_sld.tex -- MKP, SLD: subgradient iterations t* needed to beat the GNN root bound, and time of these t* iterations
% vs. one GNN bound call. Medians over 10 instances per (n, N); generated from latex_data/time_vs_tstar_sld.csv
% (extract_time_vs_tstar.py). Preamble: \usepackage{pgfplots} \usepgfplotslibrary{groupplots} \pgfplotsset{compat=1.16}
% Usage in the paper: % fig_tstar_sld.tex -- MKP, SLD: subgradient iterations t* needed to beat the GNN root bound, and time of these t* iterations
% vs. one GNN bound call. Medians over 10 instances per (n, N); generated from latex_data/time_vs_tstar_sld.csv
% (extract_time_vs_tstar.py). Preamble: \usepackage{pgfplots} \usepgfplotslibrary{groupplots} \pgfplotsset{compat=1.16}
% Usage in the paper: % fig_tstar_sld.tex -- MKP, SLD: subgradient iterations t* needed to beat the GNN root bound, and time of these t* iterations
% vs. one GNN bound call. Medians over 10 instances per (n, N); generated from latex_data/time_vs_tstar_sld.csv
% (extract_time_vs_tstar.py). Preamble: \usepackage{pgfplots} \usepgfplotslibrary{groupplots} \pgfplotsset{compat=1.16}
% Usage in the paper: \input{fig_tstar_sld}
\pgfplotstableread{
n  N     tstar  tsub      tgnn
15 100   16.5   0.03475   0.002985
15 500   20     0.1942    0.01221
15 1000  58     0.9757    0.0249
20 100   20.5   0.2625    0.01328
20 500   20.5   0.6345    0.05749
20 1000  60     2.886     0.1172
25 100   14     1.565     0.09308
25 500   13     2.407     0.4131
25 1000  33.5   6.722     0.8262
30 100   6      4.417     0.6154
30 500   9      9.976     2.424
30 1000  22     22.66     4.55
}\tstarsld
\tikzset{tstarSG/.style={line width=1pt, mark=*, mark size=1.6pt},
         tstarGNN/.style={line width=1pt, dashed, mark=square*, mark size=1.4pt}}
\definecolor{c100}{HTML}{6DA7EC}
\definecolor{c500}{HTML}{2A78D6}
\definecolor{c1000}{HTML}{104281}
\begin{figure}[t]
\centering
\begin{tikzpicture}
\begin{groupplot}[
    group style={group size=2 by 1, horizontal sep=1.45cm},
    width=0.5\linewidth, height=4.4cm, scale only axis=false,
    every axis/.append style={font=\footnotesize}, xlabel={$n$}, xtick={15,20,25,30}, xmin=13.5, xmax=31.5,
    grid=major, grid style={gray!20}, axis line style={gray!60}, tick style={gray!60},
    unbounded coords=discard,
  ]
  \nextgroupplot[ylabel={Iterations $t^\star$}, ymin=0, ymax=70,
    legend to name=legtstarN, legend columns=3, legend style={draw=none, font=\footnotesize, column sep=6pt}]
    \addplot+[tstarSG, color=c100, mark options={fill=c100}, x filter/.expression={\thisrow{N}==100 ? x : nan}] table[x=n, y=tstar] {\tstarsld};
    \addplot+[tstarSG, color=c500, mark options={fill=c500}, x filter/.expression={\thisrow{N}==500 ? x : nan}] table[x=n, y=tstar] {\tstarsld};
    \addplot+[tstarSG, color=c1000, mark options={fill=c1000}, x filter/.expression={\thisrow{N}==1000 ? x : nan}] table[x=n, y=tstar] {\tstarsld};
    \legend{$N=100$, $N=500$, $N=1000$}
  \nextgroupplot[ymode=log, ylabel={Time (s)}, ymin=1e-3, ymax=60,
    legend to name=legtstarM, legend columns=2, legend style={draw=none, font=\footnotesize, column sep=6pt}]
    \addlegendimage{tstarSG, color=black, mark options={fill=black}}
    \addlegendimage{tstarGNN, color=black, mark options={fill=black}}
    \addplot+[tstarSG, color=c100, mark options={fill=c100}, x filter/.expression={\thisrow{N}==100 ? x : nan}] table[x=n, y=tsub] {\tstarsld};
    \addplot+[tstarSG, color=c500, mark options={fill=c500}, x filter/.expression={\thisrow{N}==500 ? x : nan}] table[x=n, y=tsub] {\tstarsld};
    \addplot+[tstarSG, color=c1000, mark options={fill=c1000}, x filter/.expression={\thisrow{N}==1000 ? x : nan}] table[x=n, y=tsub] {\tstarsld};
    \addplot+[tstarGNN, color=c100, mark options={fill=c100}, x filter/.expression={\thisrow{N}==100 ? x : nan}] table[x=n, y=tgnn] {\tstarsld};
    \addplot+[tstarGNN, color=c500, mark options={fill=c500}, x filter/.expression={\thisrow{N}==500 ? x : nan}] table[x=n, y=tgnn] {\tstarsld};
    \addplot+[tstarGNN, color=c1000, mark options={fill=c1000}, x filter/.expression={\thisrow{N}==1000 ? x : nan}] table[x=n, y=tgnn] {\tstarsld};
    \legend{Subgradient ($t^\star$ iterations), GNN (one call)}
\end{groupplot}
\end{tikzpicture}\\[2pt]
\pgfplotslegendfromname{legtstarN}\\
\pgfplotslegendfromname{legtstarM}
\caption{MKP, SLD. Left: number of subgradient iterations $t^\star$ (cold start at the root) needed to reach the bound inferred by the GNN in a single forward pass. Right: wall-clock time of these $t^\star$ iterations (solid) versus one GNN bound call (dashed; graph construction, forward pass and evaluation of $L(\lambda)$, TorchScript warm-up excluded). Medians over 10 instances per $(n, N)$; 1 H100 GPU, 4 cores. For $n=15$, the subgradient time is $t^\star/t_{\max}$ times the time of a full run of $t_{\max}$ iterations.}
\label{fig:tstar_sld}
\end{figure}

\pgfplotstableread{
n  N     tstar  tsub      tgnn
15 100   16.5   0.03475   0.002985
15 500   20     0.1942    0.01221
15 1000  58     0.9757    0.0249
20 100   20.5   0.2625    0.01328
20 500   20.5   0.6345    0.05749
20 1000  60     2.886     0.1172
25 100   14     1.565     0.09308
25 500   13     2.407     0.4131
25 1000  33.5   6.722     0.8262
30 100   6      4.417     0.6154
30 500   9      9.976     2.424
30 1000  22     22.66     4.55
}\tstarsld
\tikzset{tstarSG/.style={line width=1pt, mark=*, mark size=1.6pt},
         tstarGNN/.style={line width=1pt, dashed, mark=square*, mark size=1.4pt}}
\definecolor{c100}{HTML}{6DA7EC}
\definecolor{c500}{HTML}{2A78D6}
\definecolor{c1000}{HTML}{104281}
\begin{figure}[t]
\centering
\begin{tikzpicture}
\begin{groupplot}[
    group style={group size=2 by 1, horizontal sep=1.45cm},
    width=0.5\linewidth, height=4.4cm, scale only axis=false,
    every axis/.append style={font=\footnotesize}, xlabel={$n$}, xtick={15,20,25,30}, xmin=13.5, xmax=31.5,
    grid=major, grid style={gray!20}, axis line style={gray!60}, tick style={gray!60},
    unbounded coords=discard,
  ]
  \nextgroupplot[ylabel={Iterations $t^\star$}, ymin=0, ymax=70,
    legend to name=legtstarN, legend columns=3, legend style={draw=none, font=\footnotesize, column sep=6pt}]
    \addplot+[tstarSG, color=c100, mark options={fill=c100}, x filter/.expression={\thisrow{N}==100 ? x : nan}] table[x=n, y=tstar] {\tstarsld};
    \addplot+[tstarSG, color=c500, mark options={fill=c500}, x filter/.expression={\thisrow{N}==500 ? x : nan}] table[x=n, y=tstar] {\tstarsld};
    \addplot+[tstarSG, color=c1000, mark options={fill=c1000}, x filter/.expression={\thisrow{N}==1000 ? x : nan}] table[x=n, y=tstar] {\tstarsld};
    \legend{$N=100$, $N=500$, $N=1000$}
  \nextgroupplot[ymode=log, ylabel={Time (s)}, ymin=1e-3, ymax=60,
    legend to name=legtstarM, legend columns=2, legend style={draw=none, font=\footnotesize, column sep=6pt}]
    \addlegendimage{tstarSG, color=black, mark options={fill=black}}
    \addlegendimage{tstarGNN, color=black, mark options={fill=black}}
    \addplot+[tstarSG, color=c100, mark options={fill=c100}, x filter/.expression={\thisrow{N}==100 ? x : nan}] table[x=n, y=tsub] {\tstarsld};
    \addplot+[tstarSG, color=c500, mark options={fill=c500}, x filter/.expression={\thisrow{N}==500 ? x : nan}] table[x=n, y=tsub] {\tstarsld};
    \addplot+[tstarSG, color=c1000, mark options={fill=c1000}, x filter/.expression={\thisrow{N}==1000 ? x : nan}] table[x=n, y=tsub] {\tstarsld};
    \addplot+[tstarGNN, color=c100, mark options={fill=c100}, x filter/.expression={\thisrow{N}==100 ? x : nan}] table[x=n, y=tgnn] {\tstarsld};
    \addplot+[tstarGNN, color=c500, mark options={fill=c500}, x filter/.expression={\thisrow{N}==500 ? x : nan}] table[x=n, y=tgnn] {\tstarsld};
    \addplot+[tstarGNN, color=c1000, mark options={fill=c1000}, x filter/.expression={\thisrow{N}==1000 ? x : nan}] table[x=n, y=tgnn] {\tstarsld};
    \legend{Subgradient ($t^\star$ iterations), GNN (one call)}
\end{groupplot}
\end{tikzpicture}\\[2pt]
\pgfplotslegendfromname{legtstarN}\\
\pgfplotslegendfromname{legtstarM}
\caption{MKP, SLD. Left: number of subgradient iterations $t^\star$ (cold start at the root) needed to reach the bound inferred by the GNN in a single forward pass. Right: wall-clock time of these $t^\star$ iterations (solid) versus one GNN bound call (dashed; graph construction, forward pass and evaluation of $L(\lambda)$, TorchScript warm-up excluded). Medians over 10 instances per $(n, N)$; 1 H100 GPU, 4 cores. For $n=15$, the subgradient time is $t^\star/t_{\max}$ times the time of a full run of $t_{\max}$ iterations.}
\label{fig:tstar_sld}
\end{figure}

\pgfplotstableread{
n  N     tstar  tsub      tgnn
15 100   16.5   0.03475   0.002985
15 500   20     0.1942    0.01221
15 1000  58     0.9757    0.0249
20 100   20.5   0.2625    0.01328
20 500   20.5   0.6345    0.05749
20 1000  60     2.886     0.1172
25 100   14     1.565     0.09308
25 500   13     2.407     0.4131
25 1000  33.5   6.722     0.8262
30 100   6      4.417     0.6154
30 500   9      9.976     2.424
30 1000  22     22.66     4.55
}\tstarsld
\tikzset{tstarSG/.style={line width=1pt, mark=*, mark size=1.6pt},
         tstarGNN/.style={line width=1pt, dashed, mark=square*, mark size=1.4pt}}
\definecolor{c100}{HTML}{6DA7EC}
\definecolor{c500}{HTML}{2A78D6}
\definecolor{c1000}{HTML}{104281}
\begin{figure}[t]
\centering
\begin{tikzpicture}
\begin{groupplot}[
    group style={group size=2 by 1, horizontal sep=1.45cm},
    width=0.5\linewidth, height=4.4cm, scale only axis=false,
    every axis/.append style={font=\footnotesize}, xlabel={$n$}, xtick={15,20,25,30}, xmin=13.5, xmax=31.5,
    grid=major, grid style={gray!20}, axis line style={gray!60}, tick style={gray!60},
    unbounded coords=discard,
  ]
  \nextgroupplot[ylabel={Iterations $t^\star$}, ymin=0, ymax=70,
    legend to name=legtstarN, legend columns=3, legend style={draw=none, font=\footnotesize, column sep=6pt}]
    \addplot+[tstarSG, color=c100, mark options={fill=c100}, x filter/.expression={\thisrow{N}==100 ? x : nan}] table[x=n, y=tstar] {\tstarsld};
    \addplot+[tstarSG, color=c500, mark options={fill=c500}, x filter/.expression={\thisrow{N}==500 ? x : nan}] table[x=n, y=tstar] {\tstarsld};
    \addplot+[tstarSG, color=c1000, mark options={fill=c1000}, x filter/.expression={\thisrow{N}==1000 ? x : nan}] table[x=n, y=tstar] {\tstarsld};
    \legend{$N=100$, $N=500$, $N=1000$}
  \nextgroupplot[ymode=log, ylabel={Time (s)}, ymin=1e-3, ymax=60,
    legend to name=legtstarM, legend columns=2, legend style={draw=none, font=\footnotesize, column sep=6pt}]
    \addlegendimage{tstarSG, color=black, mark options={fill=black}}
    \addlegendimage{tstarGNN, color=black, mark options={fill=black}}
    \addplot+[tstarSG, color=c100, mark options={fill=c100}, x filter/.expression={\thisrow{N}==100 ? x : nan}] table[x=n, y=tsub] {\tstarsld};
    \addplot+[tstarSG, color=c500, mark options={fill=c500}, x filter/.expression={\thisrow{N}==500 ? x : nan}] table[x=n, y=tsub] {\tstarsld};
    \addplot+[tstarSG, color=c1000, mark options={fill=c1000}, x filter/.expression={\thisrow{N}==1000 ? x : nan}] table[x=n, y=tsub] {\tstarsld};
    \addplot+[tstarGNN, color=c100, mark options={fill=c100}, x filter/.expression={\thisrow{N}==100 ? x : nan}] table[x=n, y=tgnn] {\tstarsld};
    \addplot+[tstarGNN, color=c500, mark options={fill=c500}, x filter/.expression={\thisrow{N}==500 ? x : nan}] table[x=n, y=tgnn] {\tstarsld};
    \addplot+[tstarGNN, color=c1000, mark options={fill=c1000}, x filter/.expression={\thisrow{N}==1000 ? x : nan}] table[x=n, y=tgnn] {\tstarsld};
    \legend{Subgradient ($t^\star$ iterations), GNN (one call)}
\end{groupplot}
\end{tikzpicture}\\[2pt]
\pgfplotslegendfromname{legtstarN}\\
\pgfplotslegendfromname{legtstarM}
\caption{MKP, SLD. Left: number of subgradient iterations $t^\star$ (cold start at the root) needed to reach the bound inferred by the GNN in a single forward pass. Right: wall-clock time of these $t^\star$ iterations (solid) versus one GNN bound call (dashed; graph construction, forward pass and evaluation of $L(\lambda)$, TorchScript warm-up excluded). Medians over 10 instances per $(n, N)$; 1 H100 GPU, 4 cores. For $n=15$, the subgradient time is $t^\star/t_{\max}$ times the time of a full run of $t_{\max}$ iterations.}
\label{fig:tstar_sld}
\end{figure}

\pgfplotstableread{
n  N     tstar  tsub      tgnn
15 100   16.5   0.03475   0.002985
15 500   20     0.1942    0.01221
15 1000  58     0.9757    0.0249
20 100   20.5   0.2625    0.01328
20 500   20.5   0.6345    0.05749
20 1000  60     2.886     0.1172
25 100   14     1.565     0.09308
25 500   13     2.407     0.4131
25 1000  33.5   6.722     0.8262
30 100   6      4.417     0.6154
30 500   9      9.976     2.424
30 1000  22     22.66     4.55
}\tstarsld
\tikzset{tstarSG/.style={line width=1pt, mark=*, mark size=1.6pt},
         tstarGNN/.style={line width=1pt, dashed, mark=square*, mark size=1.4pt}}
\definecolor{c100}{HTML}{6DA7EC}
\definecolor{c500}{HTML}{2A78D6}
\definecolor{c1000}{HTML}{104281}
\begin{figure}[t]
\centering
\begin{tikzpicture}
\begin{groupplot}[
    group style={group size=2 by 1, horizontal sep=1.45cm},
    width=0.5\linewidth, height=4.4cm, scale only axis=false,
    every axis/.append style={font=\footnotesize}, xlabel={$n$}, xtick={15,20,25,30}, xmin=13.5, xmax=31.5,
    grid=major, grid style={gray!20}, axis line style={gray!60}, tick style={gray!60},
    unbounded coords=discard,
  ]
  \nextgroupplot[ylabel={Iterations $t^\star$}, ymin=0, ymax=70,
    legend to name=legtstarN, legend columns=3, legend style={draw=none, font=\footnotesize, column sep=6pt}]
    \addplot+[tstarSG, color=c100, mark options={fill=c100}, x filter/.expression={\thisrow{N}==100 ? x : nan}] table[x=n, y=tstar] {\tstarsld};
    \addplot+[tstarSG, color=c500, mark options={fill=c500}, x filter/.expression={\thisrow{N}==500 ? x : nan}] table[x=n, y=tstar] {\tstarsld};
    \addplot+[tstarSG, color=c1000, mark options={fill=c1000}, x filter/.expression={\thisrow{N}==1000 ? x : nan}] table[x=n, y=tstar] {\tstarsld};
    \legend{$N=100$, $N=500$, $N=1000$}
  \nextgroupplot[ymode=log, ylabel={Time (s)}, ymin=1e-3, ymax=60,
    legend to name=legtstarM, legend columns=2, legend style={draw=none, font=\footnotesize, column sep=6pt}]
    \addlegendimage{tstarSG, color=black, mark options={fill=black}}
    \addlegendimage{tstarGNN, color=black, mark options={fill=black}}
    \addplot+[tstarSG, color=c100, mark options={fill=c100}, x filter/.expression={\thisrow{N}==100 ? x : nan}] table[x=n, y=tsub] {\tstarsld};
    \addplot+[tstarSG, color=c500, mark options={fill=c500}, x filter/.expression={\thisrow{N}==500 ? x : nan}] table[x=n, y=tsub] {\tstarsld};
    \addplot+[tstarSG, color=c1000, mark options={fill=c1000}, x filter/.expression={\thisrow{N}==1000 ? x : nan}] table[x=n, y=tsub] {\tstarsld};
    \addplot+[tstarGNN, color=c100, mark options={fill=c100}, x filter/.expression={\thisrow{N}==100 ? x : nan}] table[x=n, y=tgnn] {\tstarsld};
    \addplot+[tstarGNN, color=c500, mark options={fill=c500}, x filter/.expression={\thisrow{N}==500 ? x : nan}] table[x=n, y=tgnn] {\tstarsld};
    \addplot+[tstarGNN, color=c1000, mark options={fill=c1000}, x filter/.expression={\thisrow{N}==1000 ? x : nan}] table[x=n, y=tgnn] {\tstarsld};
    \legend{Subgradient ($t^\star$ iterations), GNN (one call)}
\end{groupplot}
\end{tikzpicture}\\[2pt]
\pgfplotslegendfromname{legtstarN}\\
\pgfplotslegendfromname{legtstarM}
\caption{MKP, SLD. Left: number of subgradient iterations $t^\star$ (cold start at the root) needed to reach the bound inferred by the GNN in a single forward pass. Right: wall-clock time of these $t^\star$ iterations (solid) versus one GNN bound call (dashed; graph construction, forward pass and evaluation of $L(\lambda)$, TorchScript warm-up excluded). Medians over 10 instances per $(n, N)$; 1 H100 GPU, 4 cores. For $n=15$, the subgradient time is $t^\star/t_{\max}$ times the time of a full run of $t_{\max}$ iterations.}
\label{fig:tstar_sld}
\end{figure}
